\section{Discussion}
\label{sec:discussion}

\subsection{Limitations}
\label{subsec:limitations}

\paragraph{Confident wrongness.} \eps cannot detect a model that is
wrong and confident. If training placed all mass on a deprecated API,
the token probability is sharp and \eps is low; the system emits the
wrong code with \textsc{complete} status. Scenario B is the
concrete example: GPT-4o emitted
\texttt{openai.ChatCompletion.create(\ldots)} --- the v0 API, which
raises \texttt{AttributeError} on SDK $\ge 1.0$ --- with no
hesitation at that token; \eps fired only because of adjacent
uncertainty in surrounding parameter design. Scenario C on
GPT-4o-mini ($\eps = 0.000$, \S\ref{subsec:scenarios}) is the cleaner
empirical instance: the model's prior on SQLAlchemy had converged
on a single version, so there is no token-level hesitation for
\eps to see. Wang et al.~\cite{wang2025deprecated} report this
failure mode at scale: $25$--$38\%$ deprecated-API emission on
standard prompts across seven LLMs, driven by confident predictions
from training data that predates the deprecation. This is a
fundamental limit of any generation-time signal and a primary
motivation for pairing \eps with static analysis
(\S\ref{subsec:static-cov}) rather than replacing it.

\paragraph{Confidently-resolved domains.} Prompts where every model
has converged on a strong training-data prior produce $\eps = 0.000$
regardless of correctness. Current-UTC datetime retrieval and
HTTP-GET-to-JSON fall here across all four models. For stable
widely-deployed APIs this is desirable; for deprecated-but-%
confidently-recalled APIs it is Scenario B's failure mode.

\paragraph{Token-level localization depends on function-level fire.}
The token-level focus story (\S\ref{subsec:token-focus}) rests on
the function being \textsc{flagged} in the first place. A confident-%
wrong function has no high-\eps tokens by definition; the token-%
level claim does not rescue the function-level miss, because there
is no miss to rescue --- the system produced \textsc{complete} and
the developer never sees a reviewer call. Token localization
sharpens the review of already-flagged functions; it does not expand
recall.

\paragraph{Small-sample calibration.} Main calibration is $30$
prompts per model; Wilson intervals on \textsc{high} detection are
$\pm 17$pp at $n=20$. We claim recall-first design plus an empirical
miss rate of zero on the sets evaluated, not a proof of perfect
recall. The \papg comparison in \S\ref{subsec:papg-comp} uses the
same $30$ prompts, which bounds all statistical claims in that
subsection to the same small-sample regime. The token-focus
analysis aggregates $144$ entries (four models, $30$ prompts with
some drop-outs) and should be read as characterising order of
magnitude rather than a precise percentage.

\paragraph{\papg comparison scope.} The signal-level \papg evaluation
(\S\ref{subsec:papg-comp}) uses the calibration benchmark ($30$
prompts per model, $120$ total). The end-to-end comparison through
the review loop (\S\ref{subsec:papg-comp}, Table~\ref{tab:pavg-endtoend})
uses the full $228$-entry production-plus-scenario set. The recall
ceiling finding --- \papg's maximum API-entry recall is $44\%$ at
any threshold --- rests on the production set and is not subject to
the small-sample caveat that bounds the calibration results.

\paragraph{Multi-sample sensitivity.} The multi-sample comparison
uses $N=5$ at $T=0.7$. Higher $N$ or higher $T$ might recover
additional behavioural diversity on some prompts --- a $T=1.0$ run
would nearly double the sampling probability of the minority branch
on a $0.52/0.47$ token --- though the structural argument does not
depend on parameter choice: \eps is read at $T=0$ where the model's
internal distribution is exposed directly, while any sampling-based
method is limited to what that distribution produces as behaviour.

\paragraph{Type 3 floor.} The Type 3 false-positive floor remains;
the mitigation is reviewer speed (both branches pass).

\paragraph{LLM-judged ground truth.}\label{subsec:gt-disclosure}
\textsc{low} ground truth is \texttt{exec()}-verified against
assertions covering return type, edge cases, and correctness.
\textsc{high} ground truth is LLM-judged via GPT-4o. Circularity
exists where the evaluated model shares a family with the judge
(GPT-4o-mini judged by GPT-4o); results for that model should be
read with that caveat. We disclose rather than overclaim human
verification we did not perform.

\paragraph{Hyperparameters.} The cascaded tier algorithm (six
checks, three thresholds) was calibrated on GPT-4o; other models
inherit these. The binary \textsc{flagged+} gate is robust;
internal tier placement is a refinement on top.

\subsection{\textsc{paused} tier and delivery semantics}
\label{subsec:paused}

\S\ref{subsec:tiers} defined \textsc{paused} as $\eps \ge 0.95$
at the function level. Whether the function-level \textsc{paused}
tier actually holds code delivery until the review completes, or
simply runs the review loop like \textsc{flagged}, is an
implementation choice we do not prescribe. Production use would
likely default to the non-blocking treatment (same routing as
\textsc{flagged}) while exposing \textsc{paused} as a higher-urgency
signal in whatever UI surfaces reviewer output.

\subsection{Context learning as a microcosm}
\label{subsec:context}

Injecting learned false-positive patterns into the reviewer's
context (the configuration evaluated in \S\ref{subsec:review-loop})
raised clearance from $53.2\%$ to $68.7\%$, and late false positives
fell from $4$ to $0$. This is a trade-off, not a free win: the same
priming that eliminates late false positives also suppresses some
genuine failures that the context-free reviewer had forwarded
(Figure~\ref{fig:review-loop}). Practitioners in high-stakes
settings should weigh the suppressed-failure cost against the
clearance and late-FP gains. This is in-context improvement on a
fixed model; training-time improvement would be expected to compound
with it. The mechanism generalizes to any setting where a downstream
model is judging upstream output with known failure patterns.

\subsection{Future work}
\label{subsec:future}

\paragraph{Per-token intervention.}
The token-level focus invites editor-level surfacing: highlight the
handful of high-\eps tokens inline as the developer reviews, the
way a spellchecker does. A targeted re-decode at lower temperature
on just the flagged span (AdaDec-style, but at review time) is a
natural follow-up.

\paragraph{Token-routed multi-sample.}
Multi-sample is structurally blind at the function level but could
recover signal if restricted to the high-\eps token span at high
temperature --- a hybrid that uses \eps to localize and sampling to
explore. At $\sim 10\%$ of tokens this costs far less than full
multi-sample.

\paragraph{Pipeline propagation.}
Upstream \eps should propagate as prior uncertainty into downstream
generations, modulating sampling temperature or review threshold.

\paragraph{Open-source models.}
vLLM exposes full top-$k$; a Llama / Qwen / Mistral calibration run
would test whether the recall-first boundary and the $10.5\%$
token-focus regularity hold outside OpenAI.

\paragraph{Cross-decision consistency.}
Scenario D's mismatches argue for AST-level consistency checks over
\eps (e.g., \texttt{async def} $+$ blocking HTTP client).

\paragraph{Import-gated activation.}
Restricting \eps scoring to functions that import external libraries
would eliminate most of the Type 3 floor at negligible cost.

